\subsection{Relative class-number selection: the exact remaining scope}
\label{KT-RELATIVE-CLASS-SELECTION-SCOPE}

This is a primary-source applicability audit, classified as
\textbf{NONMOVEMENT}.  Let $F=F_d(p)$ be the totally real subfield of
$\mathbb Q(\zeta_p)$ cut out by the even character $\omega^e$.
The proposed sufficient input is a CM compositum $FE$, with $E/\mathbb Q$
imaginary quadratic, whose relative class number is a $p$-unit.
The abelian relative class-number factorization would then package the
odd twists $\eta\omega^{je}$ together.  The present audit neither proves
that such an $E$ exists nor imports that factorization as a new theorem.

\paragraph{The 1991 source is Naito, not Horie.}
Hirotada Naito, \emph{Indivisibility of class numbers of totally imaginary
quadratic extensions and their Iwasawa invariants}, Journal of the
Mathematical Society of Japan \textbf{43} (1991), 185--194,
doi:10.2969/jmsj/04310185, states on printed p.~185 that, for totally real
$F$ and an odd prime $\ell$ satisfying
\[
             \ell\nmid w_F\zeta_F(-1),                 \tag{RC.1}
\]
there are infinitely many totally imaginary quadratic extensions $K/F$
with $\ell\nmid h(K)/h(F)$ and with every prime above $\ell$ nonsplit.
Here $w_F$ is the root-of-unity denominator defined immediately before
the theorem.  The complete construction in \S3, pp.~189--191, varies
quadratic extensions of $F$ through quaternionic trace formulas.  It
does not require their defining squareclasses to come from
$\mathbb Q^\times$.

Thus there are two distinct missing inputs for the proposed application:
the unit condition (RC.1) at $\ell=p$, and descent of the selected
quadratic extension to an imaginary quadratic field over $\mathbb Q$.
The fact that $p$ ramifies in $F_d(p)$ is not, by itself, an excluded
hypothesis of Naito's statement.  There is a further essential source
qualification: Naito's \emph{Erratum}, Journal of the Mathematical
Society of Japan \textbf{46} (1994), 725--726,
doi:10.2969/jmsj/04640725, explicitly withdraws the proof of the
proposition in \S4 of the original article.  An order needed for its
trace calculation can be omitted.  The erratum repairs particular
bounded real-quadratic numerical examples, not the general
seed-propagation proposition.  The original main theorem under
(RC.1) is unaffected.  No \S4 propagation theorem is imported here.

\paragraph{Takai's fixed-field theorem and its seed.}
Yuuki Takai, \emph{Indivisibility of relative class numbers of totally
imaginary quadratic extensions and vanishing of these relative Iwasawa
invariants}, Journal of Number Theory \textbf{186} (2018), 162--179,
doi:10.1016/j.jnt.2017.09.024, treats totally real Galois fields $F$.
Theorem~1, pp.~172--176, assumes $p>2g+1$ and
\[
 g\leq \frac{[F(\zeta_p):F]}
                 {2^{v_2([F(\zeta_p):F])}},\qquad g=[F:\mathbb Q],
                                                               \tag{RC.2}
\]
and an explicitly specified nonzero mod-$p$ finite trace-coefficient
sum.  Prescribed quadratic Hecke-character values are allowed.  Its
extensions have the form $F(\sqrt{-2\xi})$ for elements $\xi\in F$;
rational squareclasses are not imposed.  Corollary~1, pp.~176--177,
removes the coefficient premise for sufficiently large $p$ with $F$
fixed.  Theorem~2 and Remark~3, p.~178, retain a coefficient premise
when the primes above $p$ are required to be nonsplit.

The complete proof of Corollary~1 contains two precise fixed-field
steps.  First, a positive rational coefficient sum is fixed before
$p$ varies, so it becomes a $p$-unit away from a finite set.  Second,
the auxiliary field $K'$ is fixed independently of $p$, permitting
$K'\cap\mathbb Q(\zeta_p)=\mathbb Q$ for sufficiently large $p$.
Neither step gives a uniform assertion for $F=F_d(p)$.  Indeed, in
that application $F_d(p)\subset K'\cap\mathbb Q(\zeta_p)$.
Also (RC.2) becomes an actual restriction on the odd part of
$(p-1)/d$; it cannot be discarded merely because $p$ is large.
These observations do not refute Takai's theorem: they identify the
quantifiers and hypotheses that remain to be discharged here.

\paragraph{The October 2026 refinement.}
Olivia Beckwith and T\`ung Ho\`ang, \emph{Selmer groups of CM-twists
of elliptic curves}, arXiv:2609.23847v2, 8 October 2026,
Lemma~2.13 and Theorem~2.14, pp.~18--19, replace (RC.2) and the
other numerical restriction by
\[
                         [F(\zeta_p):F]\geq 2g.        \tag{RC.3}
\]
The full short proof uses that the intersection of a cyclic extension
with a multiquadratic extension has degree at most two.  Thus (RC.3)
is automatic for the quintic field $F_5(p)$ once $p\geq51$.
This refinement still retains the nonzero trace-coefficient premise
and does not restrict the output to rational squareclasses.
The elliptic-curve Theorem~4.1, pp.~25--26, additionally assumes
$p\nmid h(F)$; that application cannot remove the case where the
quintic class already survives.  Remark~4.2, p.~27, reduces the
coefficient premise to a single rational-squareclass extension when
an inert auxiliary prime $q$ with $p\nmid h^-(F(\sqrt{-q})/F)$
is already supplied.  It does not prove the existence of such a
unit seed.  The degree-at-most-five restriction appearing in the
earlier abstract is absent from this version's theorem.

\paragraph{Retrieval and checked scope.}
The official Naito PDF was retrieved, hashed, and read for its main
statement, the complete \S3 proof, and the \S4 proposition and its
withdrawn proof.  The complete two-page 1994 erratum was also read.
The trace-formula and Galois-representation foundations in \S\S1--2
remain cited inputs rather than a new foundational proof audit.
Takai's indexed publisher text supplied the complete Theorem~1 proof,
Corollary~1 proof, and the Theorem~2/Remark~3 passages.  Direct PDF
download and rendered-page retrieval returned HTTP~403.  Some indexed
formulas have visibly defective extraction; accordingly this audit
does not certify their exact coefficient normalizations or quantitative
constants.  The qualitative premises and the fixed-field steps above
are legible in the primary text.  No Takai PDF hash or full-paper
visual audit is claimed.
The official arXiv v2 PDF was retrieved and hashed.  Its exact
Lemma~2.13, Theorem~2.14, and Theorem~4.1 proofs, and Remark~4.2,
were read in full.  The preliminary twisting and elliptic-curve
foundations are cited inputs, not independently reproved here.

\paragraph{Unmet arithmetic condition.}
A result selecting general CM quadratic extensions of $F_d(p)$ is
insufficient for the source's imaginary-quadratic auxiliary field.
An applicable selection theorem must provide a rational negative
squareclass, the required local and class-number conditions, and the
simultaneous unit values at the same fixed residue prime $p$.
None of these audited statements currently provides those inputs without
an additional hypothesis.  No class component, strict Selmer group,
or $K$-group is eliminated by this audit.
